\section{Design principles and future work}

The architecture can be summarized by ten principles:

\begin{enumerate}
  \item The model is not the system of record.
  \item Durable history is authority; projections are rebuildable views.
  \item An event does not choose the state or identity context in which it is judged.
  \item Memory, context, evidence, and authority have separate trust semantics.
  \item Proposal, execution, verification, and acceptance are distinct operations.
  \item Side effects require explicit authority rather than ambient privilege.
  \item Evidence is bound to bytes, not merely names or natural-language descriptions.
  \item Changed dependencies have explicit consequences.
  \item The verifier is an attack surface and must be challenged independently.
  \item Claims about green runs are scoped to the commit and environment that produced them.
\end{enumerate}

Future work should strengthen boundaries rather than accumulate technologies: stronger external identity/attestation, more formal transition specifications and model checking, deterministic execution backends where appropriate, richer scientific verification contracts, cross-machine numerical-equivalence policies, provenance-aware context compaction, larger collusion campaigns, long-duration fault injection, performance envelopes for very large histories, schema migration rules, and additional independent readers.

A subsystem should be added because a concrete failure mode requires a boundary, not because its name makes the stack look complete.

\section{Conclusion}

The core problem is not whether an AI agent can generate scientific code or maintain a long conversation. It is whether a long-running scientific process remains reconstructible when the agent forgets, restarts, changes assumptions, produces stale evidence, encounters a different machine, or reports a state that durable history does not support.

The Scientific AI Harness makes authority external to the model. History is durable; state is replayed; memory is governed but non-authoritative; evidence resolves to bytes; dependency changes can invalidate conclusions; side effects require bounded authority; proposal and verification are separated; a second implementation can challenge the first; and the verifier itself is subjected to adversarial testing.

The intended result is not an autonomous truth machine. It is a system in which scientific automation can become more capable without being allowed to silently define its own reality.

QTA remains the originating stress case and remains forecast-only. The broader contribution is the harness around it: a governed execution and evidence substrate for long-horizon scientific agents that must act without becoming their own final authority.

\appendix
\section{Reviewer reading order}
For implementation review, a useful order is:
\begin{enumerate}
  \item \code{AGENT\_SUBSTRATE.md} -- architecture and invariants.
  \item \code{docs/DEFECT\_LEDGER.md} -- defects actually found and the claims they invalidated.
  \item \code{qta\_agent/authority.py}, \code{events.py}, and \code{store.py} -- primary authority path.
  \item \code{qta\_agent/reconstruct.py} and \code{separate\_verify.py} -- independent reading.
  \item \code{qta\_agent/memory.py} and \code{context.py} -- remembered/viewed information versus evidence.
  \item \code{qta\_agent/tasks.py}, \code{scheduler.py}, \code{netauth.py}, \code{secrets.py}, and \code{readpath.py} -- long-running work and side-effect authority.
  \item \code{qta\_agent/governed\_stage10.py} -- production integration.
  \item \code{tools/mutation\_matrix.py}, \code{tools/mutations/}, and hostile/long-horizon tests -- enforcement-point testing.
  \item \code{docs/completion\_matrix.json} -- requirement accounting.
  \item QTA scientific documentation -- case-study physics and its separate scientific claim boundary.
\end{enumerate}

\section{Snapshot discipline}
This paper intentionally avoids timeless branch-level claims such as ``complete'' or ``green.'' Run-specific evidence should be read from commit-scoped CI and \code{docs/DEFECT\_LEDGER.md}; requirement closure from \code{docs/completion\_matrix.json}; scientific gate status from QTA's scientific source of truth.

The categories remain separate:
\begin{itemize}
  \item software requirement closure is not absence of defects;
  \item absence of currently known defects is not proof of correctness;
  \item reproducible execution is not scientific validation;
  \item scientific validation is not created by the harness automatically.
\end{itemize}

That separation is part of the architecture, not a documentation convention.

